\chapter{Stability and the Yau--Tian--Donaldson bridge}
\index{Yau--Tian--Donaldson conjecture}\index{K-stability!comparison bridge}\index{entropy slope}
\label{chap:ytd-bridge}

The Yau--Tian--Donaldson bridge compares three kinds of data:
intersection numbers of degenerations, convex functionals on the
finite-energy space, and non-Archimedean functionals on valuations or
model metrics. The dictionary is exact on several large domains, but
the final existence equivalence depends on which degenerations and
which automorphism group are allowed. This chapter states every bridge
with its domain and dates the recent claims that have not been
dependency-audited for this book.

The status ledger in this chapter is frozen at the literature cutoff of
22 August 2026. A statement labelled as a dated preprint claim remains
separate from the timeless theorem chain even when its abstract announces a
stronger result. This separation is especially important here: an existence
theorem, a slope formula, and a density or regularization theorem solve
different arrows of the YTD diagram and are not interchangeable.

We begin with a smooth polarized projective manifold $(X,L)$ of
dimension $n\geq1$, and set
\[
                V=(L^n),
        \qquad
                \overline S
        =-\,\frac{n(K_X\cdot L^{n-1})}{V}.
\tag{68.1}\label{eq:ytd-volume-slope}
\]
Later, $c_1(L)$ is replaced by a K\"ahler class $\alpha$.
All potentials are considered modulo constants when the functional is
translation invariant; $E$ itself records the additive origin. Unless a
probability measure is explicitly displayed, Monge--Amp\`ere measures retain
the book's raw mass $V$. The analytic entropy used in the Mabuchi package
below is half the ordinary relative entropy, exactly as in the slope source.
The polarized automorphism group is understood modulo fibrewise scalars
on $L$, which already act by additive constants on potentials. Thus
triviality of its identity component refers to its effective action on $X$.

\section{Three functional packages}
\label{sec:ytd-three-packages}

\subsection{The analytic package}
\label{subsec:ytd-analytic}

Fix a smooth K\"ahler form $\omega\in c_1(L)$. For
$\varphi\in\mathcal E^1(X,\omega)$, put
\[
 E_\omega(\varphi)
 =\frac{1}{(n+1)V}
   \sum_{j=0}^{n}
    \int_X\varphi\,
       \omega_\varphi^j\wedge\omega^{n-j},
\qquad
 \omega_\varphi=\omega+\ddc\varphi.
\tag{68.2}\label{eq:ytd-analytic-energy}
\]
The $J$-functional is
\[
                    J_\omega(\varphi)
       =V^{-1}\int_X\varphi\,\omega^n-E_\omega(\varphi).
\tag{68.3}\label{eq:ytd-analytic-j}
\]
Mixed products here are nonpluripolar products. The energy extends by
decreasing truncation from bounded potentials, and is finite on
$\mathcal E^1$. The $J$-functional is then defined by
\eqref{eq:ytd-analytic-j}; it is not asserted to be monotone in the potential.
Both functionals are continuous for the finite-energy topology.

\begin{remark}[Energy translation]
\label{rem:ytd-energy-translation}
Adding a constant gives a useful normalization check. Since every mixed
Monge--Amp\`ere measure in \eqref{eq:ytd-analytic-energy} has mass $V$,
\[
 E_\omega(\varphi+c)=E_\omega(\varphi)+c,
 \qquad
 J_\omega(\varphi+c)=J_\omega(\varphi).
\]
Thus $E$ records the choice of additive origin, whereas $J$ descends to
potentials modulo constants. The same translation ledger will reappear for
$E^{\NA}$ and $J^{\NA}$.
\end{remark}

Let
\[
                         \mu_\varphi=V^{-1}\omega_\varphi^n
\]
and let $\mu_0=V^{-1}\omega^n$. The entropy is
\[
 H_\omega(\varphi)
  :=\frac12\operatorname{Ent}_{\mu_0}(\mu_\varphi)
  =
  \begin{cases}
   \displaystyle
   \frac12\int_X\log\!\left(\frac{\rd\mu_\varphi}{\rd\mu_0}\right)\,\rd\mu_\varphi,
      &\mu_\varphi\ll\mu_0,\\[6pt]
   +\infty,&\text{otherwise}.
  \end{cases}
\tag{68.4}\label{eq:ytd-analytic-entropy}
\]

The factor $1/2$ is \cite[equation~(1.10)]{BHJ19}. The ordinary entropy in
\Cref{prop:hybrid-entropy-liminf} is $2H_\omega$, so weak-measure entropy
results and the slope normalization must be compared with that factor.
Put $\operatorname{Ric}\omega=-\ddc\log\omega^n$, representing $c_1(X)$,
and normalize scalar curvature by
$S(\omega)\omega^n=n\operatorname{Ric}\omega\wedge\omega^{n-1}$.
For a finite-energy potential define the Ricci energy by mixed integration:
\[
 R_\omega(\varphi)
 =-\frac1{2V}\sum_{j=0}^{n-1}
   \int_X\varphi\,\operatorname{Ric}\omega\wedge
        \omega_\varphi^j\wedge\omega^{n-1-j}.
\tag{68.4a}\label{eq:ytd-analytic-ricci-energy}
\]
It is finite because the twisting form is smooth. This is the source's
$-E^{n\operatorname{Ric}\omega}$ after replacing its norm weight by
$\varphi/2$.

\begin{definition}[Mabuchi functional]
\label{def:ytd-analytic-mabuchi}
The Mabuchi functional on $\mathcal E^1(X,\omega)$ is the
lower-semicontinuous extension of the classical K-energy, normalized
by $M_\omega(0)=0$.  In the book's squared-norm convention it is
characterized on smooth potentials by
\[
 \frac{\rd}{\rd s}M_\omega(\varphi_s)
   =-\frac{1}{2V}\int_X
      \dot\varphi_s
      \bigl(S(\omega_{\varphi_s})-\overline S\bigr)
      \omega_{\varphi_s}^n.
\tag{68.5}\label{eq:ytd-mabuchi-variation}
\]
Equivalently, with the Ricci energy just defined, the Chen--Tian
decomposition is
\[
             M_\omega
       =H_\omega+\frac{\overline S}{2}E_\omega+R_\omega.
\tag{68.5a}\label{eq:ytd-chen-tian}
\]
\end{definition}

\begin{remark}[Translation of the full Mabuchi package]
\label{rem:ytd-mabuchi-translation}
Every mixed measure in the Ricci formula has integral
$V\overline S/n$ against $\operatorname{Ric}\omega$. Consequently
$R_\omega(\varphi+c)=R_\omega(\varphi)-\overline S c/2$.
The entropy is unchanged because curvature is unchanged. Its translation,
the Ricci translation and \Cref{rem:ytd-energy-translation} show directly
that $M_\omega(\varphi+c)=M_\omega(\varphi)$. Thus the origin-dependent
energy terms cancel in \eqref{eq:ytd-chen-tian}.
\end{remark}

\begin{proposition}\label{prop:ytd-analytic-convexity}
The functional $M_\omega$ is convex along finite-energy geodesics.
A smooth potential is a critical point if and only if
$\omega_\varphi$ has constant scalar curvature.
\end{proposition}

\begin{proof}
The first variation is \eqref{eq:ytd-mabuchi-variation}. Its vanishing
against all smooth zero-mean variations says that
$S(\omega_\varphi)-\overline S$ is orthogonal to every smooth
zero-mean function. Hence it is constant; its integral is zero by the
definition of $\overline S$, so it vanishes. Conversely, a cscK metric
makes the first variation zero. Convexity of the extended Mabuchi
functional along finite-energy geodesics is the imported result
\cite[Theorem~1.2]{BDL17}. That source uses
$\mathrm i\partial\bar\partial$, so its potential is
$u=\varphi/(2\pi)$ and its K-energy is $2M_\omega$ for the same curvature
forms and reference. These positive linear changes preserve finite-energy
geodesics, lower semicontinuity and convexity.
\end{proof}

\subsection{The algebraic package}
\label{subsec:ytd-algebraic}

Let $(\mathcal X,\mathcal L)$ be a normal relatively ample test
configuration for $(X,L)$, and let
$(\overline{\mathcal X},\overline{\mathcal L})$ denote its
$\mathbf P^1$-compactification.

\begin{definition}[Donaldson--Futaki invariant]
\label{def:ytd-df}
The Donaldson--Futaki invariant in the present normalization is
\[
 \DF(\mathcal X,\mathcal L)
 =V^{-1}
    \bigl(K_{\overline{\mathcal X}/\mathbf P^1}
             \cdot\overline{\mathcal L}^{\,n}\bigr)
 +\frac{\overline S}{(n+1)V}
        \bigl(\overline{\mathcal L}^{\,n+1}\bigr).
\tag{68.6}\label{eq:ytd-df}
\]
\end{definition}

The sign in the second term agrees with
\eqref{eq:ytd-volume-slope}. Replacing both $L$ and $\overline{\mathcal L}$ by their $q$th powers
leaves this normalized invariant unchanged: $V$ scales by $q^n$ and
$\overline S$ by $q^{-1}$. Normalized ramified base change need not make
$\DF$ homogeneous; the precise relation is
\Cref{prop:ytd-base-change-correction} below.

The same test configuration gives a finitely generated filtration of
\[
                          R(X,L)=\bigoplus_{m\geq0}H^0(X,mL).
\]
The total and maximal asymptotic weights define
$E^{\NA}$ and $J^{\NA}$. This filtration viewpoint is independent
of a chosen resolution.

\subsection{The non-Archimedean package}
\label{subsec:ytd-nonarchimedean}

The test configuration defines a semipositive model metric
$\phi^{\NA}$ on $L^{\An}$. Put
\[
 E^{\NA}(\phi^{\NA})
    =\frac{1}{(n+1)V}
         (\overline{\mathcal L}^{\,n+1}).
\tag{68.7}\label{eq:ytd-na-energy}
\]
If $\lambda_{\max}$ is the maximal asymptotic filtration weight, set
\[
                         J^{\NA}(\phi^{\NA})
             =\lambda_{\max}-E^{\NA}(\phi^{\NA}).
\tag{68.8}\label{eq:ytd-na-j}
\]
This is nonnegative and vanishes precisely for a translated trivial
metric on the model locus.

Translation by $c$ adds $c$ to every asymptotic filtration weight and
to $E^{\NA}$. Consequently
\[
 E^{\NA}(\phi^{\NA}+c)=E^{\NA}(\phi^{\NA})+c,
 \qquad
 J^{\NA}(\phi^{\NA}+c)=J^{\NA}(\phi^{\NA}).
\]
This matches the analytic calculation above and is the elementary check
that the two sides use the same additive quotient.

\begin{definition}[Non-Archimedean entropy and Mabuchi functional]
\label{def:ytd-na-mabuchi}
For a semipositive model metric, define
\[
 H^{\NA}(\phi^{\NA})
    =\int_{X^{\An}}A_X(v)\,
          V^{-1}\operatorname{MA}^{\NA}(\phi^{\NA}),
\tag{68.9}\label{eq:ytd-na-entropy}
\]
On a common dominating model with $p:\overline{\mathcal X}\to X$, put
$R^{\NA}=V^{-1}(p^*K_X\cdot\overline{\mathcal L}^{\,n})$.
This is \cite[Definition~7.20]{BHJ17}. Then
\[
 M^{\NA}
     =H^{\NA}+\overline S E^{\NA}+R^{\NA}.
\tag{68.10}\label{eq:ytd-na-mabuchi}
\]
Equivalently,
\[
 M^{\NA}(\mathcal X,\mathcal L)
  =\DF(\mathcal X,\mathcal L)
   +V^{-1}
     \bigl((\mathcal X_{0,\mathrm{red}}-\mathcal X_0)
               \cdot\mathcal L^n\bigr).
\tag{68.11}\label{eq:ytd-mabuchi-df}
\]
\end{definition}

\begin{remark}[Why the central-fibre correction is present]
\label{rem:ytd-central-correction}
If $\mathcal X_0$ is nonreduced, the discrepancy integral counts
the reduced divisorial valuations, while the intersection
$\DF$ retains scheme multiplicities. The last term of
\eqref{eq:ytd-mabuchi-df} reconciles them. Thus
$M^{\NA}=\DF$ for reduced central fibre, not in general.
\end{remark}

\begin{proposition}[Normalized base change and the correction]
\label{prop:ytd-base-change-correction}
Write $\mathcal X_0=\sum_E b_E E$ and set
\[
 C(\phi)=V^{-1}\sum_E(1-b_E)(E\cdot\mathcal L^n)\leq0.
\tag{68.11a}\label{eq:ytd-correction-sum}
\]
If $\phi_d$ is the metric of the normalized degree-$d$ base change, then
\[
 M^{\NA}(\phi_d)=dM^{\NA}(\phi),\qquad
 \DF(\phi_d)=d\DF(\phi)+dC(\phi)-C(\phi_d).
\tag{68.11b}\label{eq:ytd-base-change-df}
\]
For sufficiently divisible $d$, the new central fibre is reduced and
$\DF(\phi_d)=dM^{\NA}(\phi)\leq d\DF(\phi)$.
\end{proposition}

\begin{proof}
Relative ampleness makes every component degree nonnegative, proving the
sign of $C$. The log-relative-canonical formula and projection formula give
homogeneity of $M^{\NA}$; the exact imported interface is
\cite[Proposition~7.14]{BHJ17}. Subtracting the two instances of
\eqref{eq:ytd-mabuchi-df} proves \eqref{eq:ytd-base-change-df}.
Reduction for sufficiently divisible degree is
\cite[Proposition~7.16]{BHJ17}; then $C(\phi_d)=0$ and the last assertion
follows. The function action is $D_d\phi(v)=d\phi(v/d)$ of
\Cref{ex:functor-base-change-action}, with the valuation unit stated there.
\end{proof}

\begin{proposition}\label{prop:ytd-birational-invariance}
The functionals $E^{\NA},J^{\NA},H^{\NA},R^{\NA}$, and
$M^{\NA}$ are unchanged after pulling a normal test configuration to
a dominating normal model.
\end{proposition}

\begin{proof}
Energy and the Ricci term are intersection numbers, so the projection
formula applies. The Monge--Amp\`ere measure pushes forward under the
model map; log discrepancy is evaluated at the same divisorial
valuation, so its integral is unchanged. The formula for $J^{\NA}$
uses the same asymptotic weights after integral closure of the Rees
algebra. Equation \eqref{eq:ytd-na-mabuchi} gives the last assertion.
\end{proof}

\section{Slope formulae and their domains}
\label{sec:ytd-slopes}

Let $(\varphi_s)_{s\geq0}$ be a ray associated to
$(\mathcal X,\mathcal L)$, with $s=-\log|t|$.
The variable $\varphi_s$ is the book's squared-norm potential.  If
$\psi_s=-\log|\cdot|$ denotes the ordinary norm weight used in the
complex-to-non-Archimedean slope literature, then
\[
                         \varphi_s=2\psi_s.
\]
Since the two $\ddc$-conventions give the same curvature form,
\[
 E_\omega(\varphi_s)=2E_{\omega,\mathrm{wt}}(\psi_s),
 \qquad
 J_\omega(\varphi_s)=2J_{\omega,\mathrm{wt}}(\psi_s).
\]
Accordingly every $E$- or $J$-slope below that is compared with the
standard $E^{\NA}$ or $J^{\NA}$ has denominator $2s$.  Bundled
statements quoted directly from \cite{BHJ19} will keep the source's
norm-weight notation visible.

The factor two does not multiply every functional. With the curvature form
held fixed, the Monge--Amp\`ere probability measure is the same for
$\varphi_s$ and $\psi_s$, so
\[
 H_\omega(\varphi_s)=H_{\omega,\mathrm{wt}}(\psi_s).
\]
The coefficient $1/2$ in \eqref{eq:ytd-mabuchi-variation} similarly makes
the book's Mabuchi value equal to the norm-weight Mabuchi value, after the
common additive normalization. Thus $E$ and $J$ use $2s$, whereas
entropy, Ricci energy, and Mabuchi use $s$. Mixing the two denominators
would create a spurious factor two in the Chen--Tian decomposition.

\begin{theorem}\label{thm:ytd-energy-slopes}
For the canonical weak geodesic ray of a normal ample test
configuration,
\[
 \lim_{s\to+\infty}\frac{E_\omega(\varphi_s)}{2s}
       =E^{\NA}(\phi^{\NA}),
 \qquad
 \lim_{s\to+\infty}\frac{J_\omega(\varphi_s)}{2s}
       =J^{\NA}(\phi^{\NA}).
\tag{68.12}\label{eq:ytd-energy-slopes}
\]
\end{theorem}

\begin{proof}
For the norm-weight ray $\psi_s$,
\Cref{thm:rf-energy-slope} gives
\[
       \lim_{s\to+\infty}\frac{E_{\mathrm{wt}}(\psi_s)}s
          =E^{\NA}(\phi^{\NA}).
\]
That theorem was proved from the RWN energy identity and the BHJ
intersection-moment formula, so no Deligne-pairing argument is repeated
here. Since $\varphi_s=2\psi_s$, the analytic energy in the squared-norm
convention is twice the norm-weight energy. Division by $2s$ proves the
first equality in \eqref{eq:ytd-energy-slopes}.

The linear term in $J$ is the maximal asymptotic weight; its radial slope
is $\lambda_{\max}$. Subtracting the energy slope gives
$\lambda_{\max}-E^{\NA}=J^{\NA}$, proving the second equality. This is
equivalently the hybrid intersection computation of
\Cref{prop:hybrid-energy-slope}. Neither step uses entropy regularization.
\end{proof}

\begin{theorem}\label{thm:ytd-smooth-functional-slopes}
Assume that an ample test configuration carries a smooth strictly
positive $S^1$-invariant metric near its central fibre, and let
$(\varphi_s)$ be the corresponding smooth positive ray in the book's
squared-norm convention.  Put $\psi_s=\varphi_s/2$.  Then
\[
 \lim_{s\to+\infty}\frac{F_{\mathrm{wt}}(\psi_s)}s
                  =F^{\NA}(\phi^{\NA})
\]
for
\[
                       F\in\{E,H,I,J,M,R\}.
\]
Here $F_{\mathrm{wt}}$ denotes exactly the source's norm-weight
normalization.  In particular, its $E$- and $J$-components translate to
\[
 \lim_{s\to+\infty}\frac{E_\omega(\varphi_s)}{2s}
       =E^{\NA}(\phi^{\NA}),
 \qquad
 \lim_{s\to+\infty}\frac{J_\omega(\varphi_s)}{2s}
       =J^{\NA}(\phi^{\NA}).
\tag{68.13}\label{eq:ytd-smooth-slopes}
\]
\end{theorem}

\begin{proof}
Apply \Cref{thm:rf-bhj-unaffected-slopes} to the norm-weight ray
$\psi_s$. Its expanded proof uses only the Deligne-pairing and snc
integral estimates from Sections~1--4 of \cite{BHJ19}, whose validity is
explicitly preserved by \cite{BHJ19err}. It gives
\[
 \lim_{s\to+\infty}\frac{F_{\mathrm{wt}}(\psi_s)}s
        =F^{\NA}(\phi^{\NA})
 \quad\text{for }F\in\{E,H,I,J,M,R\}.
\]
This is the first assertion, with no new slope input.

It remains only to translate conventions. In the book's squared-norm
normalization $\varphi_s=2\psi_s$. The factor-two ledger in
\eqref{eq:rf-factor-two-slope-quotients} gives
\[
 \frac{E_\omega(\varphi_s)}{2s}
       =\frac{E_{\mathrm{wt}}(\psi_s)}s,
 \qquad
 \frac{J_\omega(\varphi_s)}{2s}
       =\frac{J_{\mathrm{wt}}(\psi_s)}s.
\]
Taking limits proves \eqref{eq:ytd-smooth-slopes}. No claim about a weak
finite-energy ray is made in this normalization step.
\end{proof}

The smooth ray in
\Cref{thm:ytd-smooth-functional-slopes} is not automatically
the canonical weak geodesic ray. For weak finite-energy rays, the
following assembly statement isolates the additional entropy-slope input
that is required.

\begin{proposition}\label{prop:ytd-mabuchi-liminf}
Let $(\varphi_s)$ be a weak finite-energy ray with associated
non-Archimedean metric $\phi^{\NA}$, and denote by $R_\omega$ the
Ricci-energy term normalized so that, in the book's squared-potential
convention, the Chen--Tian formula reads
\[
 M_\omega=H_\omega+\frac{\overline S}{2}E_\omega+R_\omega.
\]
Assume that the energy and
Ricci-energy terms in the Chen--Tian decomposition satisfy the following
finite real limits:
\[
 \lim_{s\to+\infty}\frac{E_\omega(\varphi_s)}{2s}
      =E^{\NA}(\phi^{\NA}),
 \qquad
 \lim_{s\to+\infty}\frac{R_\omega(\varphi_s)}s
      =R^{\NA}(\phi^{\NA}),
\]
and, separately, that
\[
 H^{\NA}(\phi^{\NA})
       \leq
 \liminf_{s\to+\infty}\frac{H_\omega(\varphi_s)}s.
\]
Then
\[
       M^{\NA}(\phi^{\NA})
          \leq
       \liminf_{s\to+\infty}\frac{M_\omega(\varphi_s)}s
\tag{68.14}\label{eq:ytd-mabuchi-liminf}
\]
\end{proposition}

\begin{proof}
Divide the Chen--Tian decomposition by $s$. Add the assumed limits of
the energy terms to the displayed entropy inequality, and use
\eqref{eq:ytd-na-mabuchi}. Both added limits are finite; the entropy
inequality may have value $+\infty$, in which case the conclusion is
understood in the extended-real sense. No difference of infinite quantities
is taken.
\end{proof}

\begin{remark}[Why weak convergence alone is insufficient]
\label{rem:ytd-entropy-slope-firewall}
\Cref{prop:hybrid-entropy-liminf} concerns ordinary relative
entropy for weakly converging probability measures. It does not by
itself identify the rescaled quantity $s^{-1}H_\omega(\varphi_s)$
with the non-Archimedean entropy. To verify the entropy-slope hypothesis
above by a hybrid argument, one must also construct the relevant varying
reference measures, control their
logarithmic densities at radial scale, and identify the central
functional with $H^{\NA}$. No unconditional weak-ray Mabuchi
inequality is deduced here without that separate theorem.
\end{remark}

\section{The BHJ erratum quarantine}
\label{sec:ytd-erratum}

\begin{theorem}\label{thm:ytd-erratum-quarantine}
The following statements of \cite{BHJ19} are not used as established
results in this book:
\[
 \text{Theorem~5.6},\quad
 \text{Theorem~5.4},\quad
 \text{Theorem~C},\quad
 \text{Corollaries~D--E}.
\tag{68.15}\label{eq:ytd-quarantined-results}
\]
Theorem~A, Corollary~B, and all results through Section~4 remain
available.
\end{theorem}

\begin{proof}
The published erratum \cite{BHJ19err} identifies an incorrect
orbit-closure assertion in the proof of Theorem~5.6 and says that the
assertion is only possibly true. It lists Theorem~5.4, Theorem~C, and
Corollaries~D--E as dependent consequences. The same erratum explicitly
confirms Theorem~A, Corollary~B, and Sections~1--4.
\end{proof}

\begin{remark}[No laundering through terminology]
\label{rem:ytd-no-laundering}
We do not infer coercivity on all rays from a statement about smooth
test-configuration rays, and we do not call a Bergman-level assertion
a finite-energy theorem. Any post-erratum implication below comes from
an independent source.
\end{remark}

\section{Valuative criteria}
\label{sec:ytd-valuative}

In a projective Fano setting take $L=-K_X$. The familiar invariants of a prime
divisor $F$ over $X$ are
\[
 \beta(F)
  =A_X(F)V-\int_0^{+\infty}\vol(L-xF)\,\rd x,
\tag{68.16}\label{eq:ytd-beta-divisor}
\]
and
\[
  j(F)
  =\int_0^{\tau(F)}
       \bigl(V-\vol(L-xF)\bigr)\,\rd x
\tag{68.17}\label{eq:ytd-j-divisor}
\]
where $\tau(F)=\sup\{x>0:\vol(L-xF)>0\}$ is the pseudoeffective
threshold.  The cut-off must be \emph{imposed} in the second formula, not
merely observed: past $\tau(F)$ the integrand of
\eqref{eq:ytd-j-divisor} equals $V>0$, so extending the range to
$(0,+\infty)$ would diverge.  In \eqref{eq:ytd-beta-divisor} the integrand
does vanish there, and that integral may be written over
$(0,+\infty)$. These formulas require the Fano normalization when used as
a Ding criterion.

\begin{example}[The Fano projective line]
\label{ex:ytd-fano-point}
Let $X=\mathbf P^1$, $L=-K_X=\mathcal O(2)$ and $F=\{p\}$.
Then $V=2$, $A_X(F)=1$, $\tau(F)=2$ and
$\vol(L-xF)=\max\{2-x,0\}$. Hence
\[
 \beta(F)=2-\int_0^2(2-x)\,\rd x=0,
 \qquad j(F)=\int_0^2 x\,\rd x=2.
\tag{68.17a}\label{eq:ytd-fano-point-invariants}
\]
Extending the second integral beyond its threshold adds $2(T-2)$ up to
$T>2$ and therefore gives $+\infty$ as $T\to\infty$. The finite cut-off
is part of the invariant. This cscK variety has continuous automorphisms,
so it is outside the finite-automorphism hypothesis of
\Cref{thm:ytd-fano-bridge}.
\end{example}

For a general K\"ahler class $\alpha$, the audited
measure-theoretic criterion is better suited to model stability.
Let
\[
                \mu=\sum_{i=1}^Na_i\delta_{v_i},
       \qquad a_i\geq0,\quad\sum_i a_i=1,
\tag{68.18}\label{eq:ytd-divisorial-measure}
\]
be a divisorial probability measure. Here $\Ent$ is the reduced
non-Archimedean entropy of \Cref{chap:mpwn-k-stability}; the total-space
discrepancy shift of \eqref{eq:mp-total-reduced-mabuchi} is not included.
The class derivative keeps $\mu$ fixed. Define
\[
             \beta_\alpha(\mu)
            :=\Ent(\mu)+\nabla_{K_X}E_\alpha^\vee(\mu).
\tag{68.19}\label{eq:ytd-beta-measure}
\]

\begin{theorem}\label{thm:ytd-valuative-measures}
The normalized model inequality of
\Cref{def:mpwn-normalized-model-test} is equivalent to the existence of
$\delta>0$ such that
\[
            \beta_\alpha(\mu)
               \geq\delta E_\alpha^\vee(\mu)
\tag{68.20}\label{eq:ytd-valuative-criterion}
\]
for every divisorial probability measure $\mu$ of
\eqref{eq:ytd-divisorial-measure}.
Stability for all big models implies this condition.
\end{theorem}

\begin{proof}
The non-Archimedean Calabi--Yau theorem identifies a finite-energy
potential with its Monge--Amp\`ere measure. Under this identification,
\[
      \beta_\alpha(\operatorname{MA}_\alpha(\phi))
             =M_\alpha^{\NA}(\phi),
      \qquad
      E_\alpha^\vee(\operatorname{MA}_\alpha(\phi))
             \asymp J_\alpha(\phi).
\tag{68.21}\label{eq:ytd-beta-mabuchi}
\]
A divisorial measure has a PL-envelope solution, and the
Monge--Amp\`ere measure of a PL envelope is divisorial. Thus
\eqref{eq:ytd-valuative-criterion} is exactly the
normalized model inequality transported through the
Calabi--Yau correspondence, up to the dimension-dependent
comparison of $E_\alpha^\vee$ and $J_\alpha$.
\Cref{thm:mpwn-valuative-criterion} gives the details using
\cite[Theorems~10.0.1 and~10.0.4]{MPWN25}.
For the distinction from the literal all-big-model condition,
see \Cref{rem:mpwn-unnormalized-model-boundary}.
\end{proof}

\begin{remark}[Not merely one divisor at a time]
\label{warn:ytd-valuative-measures}
The criterion tests finite weighted sums of divisorial Dirac masses.
Checking only $\delta_v$ is a different condition. No dreaminess
assumption is part of
\Cref{thm:ytd-valuative-measures}.
\end{remark}

\section{Coercivity and automorphisms}
\label{sec:ytd-coercivity}

\begin{definition}[Analytic coercivity]
\label{def:ytd-coercivity}
The Mabuchi functional is $J$-coercive if there are
$\delta>0$ and $C\in\mathbf R$ such that
\[
                   M_\omega(\varphi)
          \geq\delta J_\omega(\varphi)-C
\tag{68.22}\label{eq:ytd-analytic-coercivity}
\]
for every $\varphi\in\mathcal E^1(X,\omega)$.
\end{definition}

\begin{definition}[Model uniform stability]
\label{def:ytd-model-uniform}
A compact K\"ahler pair $(X,\alpha)$ is uniformly K-stable for
models if there is $\delta>0$ such that
\[
         M^{\NA}(\mathcal X,\mathcal A)
              \geq\delta J^{\NA}(\mathcal X,\mathcal A)
\tag{68.23}\label{eq:ytd-model-uniform}
\]
for every \emph{big} analytic test configuration in the precise sense
of \cite[Definition~1.1.2]{MPWN25}, using the raw intersection
expressions \eqref{eq:mpwn-model-mabuchi} and \eqref{eq:mpwn-model-j}.
This is \Cref{def:mpwn-two-stabilities}; it implies the normalized
model inequality, whose equivalence with the full raw condition
is not used here.
\end{definition}

Testing big models is stronger than testing only K\"ahler test
configurations. The distinction is essential: the completion argument
passes through envelopes of big vertical classes.

\begin{proposition}\label{prop:ytd-coercivity-minimizer}
Assume \eqref{eq:ytd-analytic-coercivity}. Assume moreover that every
normalized $J_\omega$-bounded minimizing sequence has a
$d_1$-convergent subsequence, that $M_\omega$ is lower semicontinuous
along this convergence, and that every finite-energy minimizer of
$M_\omega$ in the class is smooth. Then $M_\omega$ attains its minimum,
and its minimizer is a smooth cscK metric.
\end{proposition}

\begin{proof}
Take a minimizing sequence $(\varphi_j)$ with a fixed additive
normalization. Coercivity gives
\[
 \delta J_\omega(\varphi_j)
       \leq M_\omega(\varphi_j)+C,
\]
so $J_\omega(\varphi_j)$ is bounded. The compactness hypothesis supplies
a $d_1$-convergent subsequence with limit $\varphi_\infty$. Lower
semicontinuity gives
\[
 M_\omega(\varphi_\infty)
 \leq\liminf_jM_\omega(\varphi_j)
 =\inf_{\mathcal E^1}M_\omega,
\]
so $\varphi_\infty$ minimizes $M_\omega$. The stated regularity hypothesis makes $\varphi_\infty$ smooth. At a
smooth minimizer every smooth first variation vanishes, and then
\eqref{eq:ytd-mabuchi-variation} shows that its scalar curvature equals
$\overline S$. This proof displays separately the coercivity,
compactness, lower-semicontinuity, and regularity inputs.
\end{proof}

\begin{remark}[Automorphism boundary]
\label{rem:ytd-automorphisms}
If $\Aut^0(X,L)\neq\{1\}$, ordinary $J$-coercivity cannot hold
along all automorphism orbits. One replaces $J$ by a reduced norm
and uses relative stability. Statements below that are written with
ordinary $J$ belong to the unrelative coercive scope, typically with
discrete automorphism group.
\end{remark}

\section{Audited existence implications}
\label{sec:ytd-audited-implications}

\begin{theorem}\label{thm:ytd-fano-bridge}
Let $X$ be a smooth projective Fano variety with finite automorphism
group.  Write $D$ and $D^{\NA}$ for the Archimedean and
non-Archimedean Ding functionals in the normalization of \cite{BBJ21}.
The following are equivalent:
\begin{enumerate}
\item there are $\delta,C>0$ such that
      \[
             D(\varphi)\geq\delta J(\varphi)-C
             \quad\text{on }\mathcal E^1(X,-K_X);
      \]
\item there is $\delta'>0$ such that
      \[
             D^{\NA}(\phi)\geq\delta'J^{\NA}(\phi)
             \quad\text{for every semipositive model metric }\phi;
      \]
\item $c_1(-K_X)$ contains a K\"ahler--Einstein metric.
\end{enumerate}
\end{theorem}

\begin{proof}
We indicate which implications are slope calculations and which are deep
variational inputs.

Assume first that (1) holds, and initially let $\phi$ be represented by
an ample test configuration carrying a smooth positive invariant metric.
Its norm weights define a smooth psh ray $(\psi_s)$.  The Ding slope
formula of \cite[Theorem~4.3]{BHJ19} and the $J$-slope formula already
audited in \Cref{thm:rf-bhj-unaffected-slopes} give
\[
 \lim_{s\to+\infty}\frac{D(\psi_s)}s=D^{\NA}(\phi),
 \qquad
 \lim_{s\to+\infty}\frac{J(\psi_s)}s=J^{\NA}(\phi).
\]
Divide the coercive inequality by $s$ and let $s\to+\infty$.  This
proves the inequality, with $\delta'=\delta$, for ample model metrics.

For a general semipositive model metric, pass to a common dominating model
on which a fixed relatively ample model metric $\phi_A$ is determined,
and set
$\phi_\varepsilon=(1-\varepsilon)\phi+\varepsilon\phi_A$
with rational $0<\varepsilon<1$.  A positive rational combination of a
relatively nef model line bundle and a relatively ample one is relatively
ample, so $\phi_\varepsilon$ is represented by an ample model and
converges uniformly to $\phi$.  The energy and
$J^{\NA}$ are continuous under this convergence.  Moreover
\[
 L^{\NA}(\phi)=\inf_v\bigl(A_X(v)+\phi(v)\bigr)
\]
is $1$-Lipschitz for the uniform norm; hence
$D^{\NA}=L^{\NA}-E^{\NA}$ is continuous as well.  Letting
$\varepsilon\downarrow0$ proves (2) for every semipositive model metric.
The cited Ding slope theorem lies in Section~4 and is outside the
quarantined part of \cite{BHJ19}.

The converse implication (2)$\Rightarrow$(1) is the main compactness
theorem \cite[Theorem~5.1]{BBJ21}.  Its proof is a foundational imported
input.  Its architecture is as follows.  If coercivity fails, normalized
finite-energy geodesic segments are extracted, using $d_1$-compactness,
to a nonconstant geodesic ray.  The nonconstant normalized ray has energy $E(U_s)=cs$ with $c<0$.
The source constructs semipositive model approximants $\phi_m$ with
$\sup\phi_m=0$, $L^{\NA}(\phi_m)\to L^{\NA}(U^{\NA})$ and
$E^{\NA}(\phi_m)\geq c$; it does not require equality of the energy of
an arbitrary ray with that of its non-Archimedean datum. After decreasing
the stability constant to $0<\varepsilon<1$, (2) implies
$L^{\NA}(\phi_m)\geq(1-\varepsilon)E^{\NA}(\phi_m)$.
The ray estimate gives $L^{\NA}(U^{\NA})\leq c$. Passing to the limit
would yield $c\geq(1-\varepsilon)c$, impossible for $c<0$.
These are \cite[Lemmas~5.7--5.8 and the proof of Theorem~5.1]{BBJ21}.  The compactness, model approximation, and singularity
comparison are the deep content of the cited theorem; they are not
silently inferred from the smooth slope formula above.

Finally, the finite-energy variational package
\cite[Lemma~2.8, Theorem~2.19, Lemma~2.20, and
Corollary~5.2]{BBJ21} identifies coercivity of $D$, in the present
finite-automorphism setting, with existence and uniqueness of its minimizer.
The Euler--Lagrange
equation of a minimizer is the K\"ahler--Einstein equation, and the
regularity theorem in the smooth Fano case upgrades the minimizing current
to a smooth K\"ahler--Einstein metric.  Conversely, uniqueness of the
K\"ahler--Einstein metric implies Ding coercivity by the same cited
variational theorem.  Since $\Aut(X)$ is finite, Bando--Mabuchi
uniqueness modulo $\Aut^0(X)$ is ordinary uniqueness.  This proves the
equivalence of (1)--(3).

When $\Aut^0(X)\neq\{1\}$, the ordinary inequalities in (1) and (2)
fail along automorphism directions.  A reduced or relative theorem requires
additional definitions and inputs and is not asserted by this statement.
\end{proof}

\begin{theorem}\label{thm:ytd-mpwn-implication}
As proved in version~1 of \cite{MPWN25}, submitted
11 September 2025, if a compact K\"ahler pair $(X,\alpha)$ is
uniformly K-stable for \emph{models} in the sense of
\Cref{def:ytd-model-uniform}, then the associated
finite-energy non-Archimedean Mabuchi functional is coercive and the
existence theorem used there yields a cscK metric in $\alpha$.
In the unrelative discrete-automorphism scope, the metric is unique.
\end{theorem}

\begin{proof}
The hypothesis for all big models applies in particular to models
with $f_D\geq1$.  By \Cref{thm:mpwn-model-e1-equivalence}, this normalized
model inequality implies coercivity on the completed finite-energy
non-Archimedean space.  The proof uses entropy-preserving approximation
of finite-energy measures by measures on dual complexes and the
non-Archimedean Calabi--Yau homeomorphism. The independent existence
implication \cite[Theorem~A]{MP24} then produces a cscK metric.
This is exactly the chain stated in
\cite[Corollary~9.2.2]{MPWN25}; it does not pass through any
quarantined result of \cite{BHJ19}.
\end{proof}

\begin{remark}[Exact scope]
\label{rem:ytd-mpwn-scope}
\Cref{thm:ytd-mpwn-implication} begins with stability for all
big models. It does not prove that ordinary uniform stability tested
only on relatively K\"ahler or algebraic test configurations implies
model uniform stability. It is therefore not the unrestricted ordinary
YTD equivalence.
\end{remark}

\section{Maximal destabilizers and conditional quantization}
\label{sec:ytd-maximal-destabilizers}

Fix a smooth complex projective variety $X$ and an ample line bundle
$L$. A different use of quantization asks whether the most destabilizing
finite-dimensional direction approaches the most destabilizing geodesic
ray. The finite-dimensional minimization and the passage to the limit
have different inputs. The following account follows the known results
and explicit assumptions of \cite{Yao25MaximalDestabilizers}, version one,
submitted 20 November 2025.

\subsection{The finite-dimensional direction}

For a fixed embedding $X\hookrightarrow\mathbb P(V^*)$, put
$N=\dim V$. An additive non-Archimedean norm $\chi$ on $V$ has jumping
numbers $\lambda_1,\ldots,\lambda_N$, counted with multiplicity. Use
\[
 E_V(\chi)=\frac1N\sum_i\lambda_i,
 \qquad
 \|\chi\|_2^2=\frac1N\sum_i\lambda_i^2.
\]
The Chow weight in \cite[Theorem~33]{Yao25MaximalDestabilizers} is
\[
 M_X(\chi)=E_{\mathcal O(1)|X}(\operatorname{FS}(\chi))-E_V(\chi).
\]
Here the energy is normalized by the volume of the indicated
polarization. If the embedding is defined by $kL$, and
$\operatorname{FS}_k$ includes division of weights by $k$, this becomes
\[
 M_{X,k}(\chi)=kE^{\NA}(\operatorname{FS}_k(\chi))-E_V(\chi).
\]
In particular, the factor $k$ belongs to the energy term.

The known \cite[Theorem~20]{Yao25MaximalDestabilizers} gives a minimizer
of $M_X(\chi)/\|\chi\|_p$ for $1\leq p\leq\infty$; if $X$ is
Chow-unstable, it is unique up to positive rescaling for $1<p<\infty$.
For $p=2$ this is the maximal Chow-destabilizer. Existence uses a
fixed Hermitian reference metric: self-adjoint initial velocities
with unit spectral norm form a compact sphere in an ordinary
finite-dimensional vector space. The Chow slope is the supremum of
continuous finite-time slope quotients and is lower semicontinuous
there. This argument does not require local compactness of the space
of non-Archimedean norms in its own metric.

For uniqueness, two unit minimizers can be diagonalized in a common
basis. Convexity bounds the Chow weight of their joining segment by
the common negative minimum. Unless the endpoints agree, strict
convexity of the spectral $p$-norm makes the interior norm smaller
than one; dividing the negative weight by that smaller norm contradicts
minimality. Negativity is essential. In the semistable case scalar
norms already have weight zero; the centered formulation in
\cite[Remark~21]{Yao25MaximalDestabilizers} removes this degeneracy.

If $G\to\operatorname{GL}(V)$ preserves the embedded variety, both
the Chow weight and the spectral norm are $G$-invariant. Uniqueness
and preservation of the norm imply that the maximal destabilizer
itself is $G$-fixed, by \cite[Theorem~23]{Yao25MaximalDestabilizers}.
The known \cite[Theorem~36]{Yao25MaximalDestabilizers} also gives
\[
 E_V(\chi)=0,
 \qquad
 \operatorname{SN}_k\operatorname{FS}_k(\chi)=\chi
\]
for a maximal Chow-destabilizer of the $kL$ embedding. Here
$\operatorname{SN}_k$ is the additive section sup-norm operator.
The second identity is a saturation statement, stronger than
invariance under a group.

\subsection{The limit problem and its assumptions}

In this subsection $\|\phi\|_2$ denotes the radial $L^2$ speed in
Yao's convention $t=-\log|\tau|^2$. Thus $t=2s$ for the clock
$s=-\log|\tau|$ used earlier in this chapter, and this speed is
half the speed of the corresponding squared-potential ray when
parametrized by $s$. The source's analytic K-energy is twice the
book's $M_\omega$ at the same squared potential. Dividing their
slopes by $t$ and $s$, respectively, gives the same $M^{\NA}$.
All spectral and radial norms below retain the source's $t$ units.

For a geodesically unstable polarized manifold, the maximal
destabilizing $L^2$ ray used in
\cite[Theorem~37 and Definition~39]{Yao25MaximalDestabilizers}
determines a non-Archimedean direction $\phi^K$. The identification
of its radial Mabuchi slope with $M^{\NA}$ is assumption A1 in that
paper. In the present smooth complex projective ample scope this
is a known source result: \cite[Theorem~C]{BJ25WeightedYTD}, or
\cite[Corollary~8.4 and equation~(14)]{DarvasZhang25YTD}, gives
the equality on maximal finite-energy rays. In the first source
the entropy input is Theorem~7.5; Theorem~8.7 gives its weighted
extension. These are identified imports in this scope; the
selective source checks do not certify every antecedent of the
general weighted correspondence. They do not extend the equality
to arbitrary singular, nonprojective or nonample settings.

With finite $L^2$ speed and negative $M^{\NA}$, the homogeneous
ratio $M^{\NA}(\phi)/\|\phi\|_2$ has the canonical representative
\[
 \phi^\flat=
 \frac{-M^{\NA}(\phi^K)}{\|\phi^K\|_2^2}*\phi^K.
\]
Here $a*\phi$ rescales the ray's time by $a>0$. As explained in
\cite[Proposition~44]{Yao25MaximalDestabilizers}, the quadratic
penalty selects this representative:
\[
 \widetilde M(\phi)=M^{\NA}(\phi)+\frac12\|\phi\|_2^2,
 \qquad
 \|\phi^\flat\|_2^2=-M^{\NA}(\phi^\flat).
\]
Indeed, on any negative direction the minimizing scale of
$aM^{\NA}(\phi)+a^2\|\phi\|_2^2/2$ is
$-M^{\NA}(\phi)/\|\phi\|_2^2$.

The quantization argument retains three further assumptions from
\cite[\S7.1.2]{Yao25MaximalDestabilizers}:
\begin{enumerate}
\item[A2.] The maximal non-Archimedean direction is continuous.
\item[A3.] For every continuous non-Archimedean psh metric $\phi$,
\[
 E_k(\operatorname{SN}_k\phi)
 =E^{\NA}(\phi)-\frac{M^{\NA}(\phi)}{2k}+o(k^{-1}),
 \qquad E_k(\chi)=\frac{E_V(\chi)}k.
\]
\item[A4.] The minimizers of the approximating penalized functionals
converge to the minimizer of $\widetilde M$.
\end{enumerate}
The known \cite[Theorem~46]{Yao25MaximalDestabilizers} proves A3,
with remainder $O(k^{-2})$, for metrics of ample test
configurations along sufficiently divisible degrees. Its proof
uses reduced central fibre after normalized ramified base change,
the identification of section sup-norms with the weight filtration,
and the two-term weight expansion. This does not prove A3 for
every continuous metric. Independently,
\cite[Theorem~45]{Yao25MaximalDestabilizers} gives
\[
 \|\phi\|_2=\lim_{k\to\infty}k^{-1}
                 \|\operatorname{SN}_k\phi\|_2
\]
for continuous non-Archimedean psh metrics.

\subsection{What the reduction establishes}

Set
\[
 Q_k(\phi)=2k\bigl(E^{\NA}(\phi)-E_k(\operatorname{SN}_k\phi)\bigr),
 \qquad
 \widetilde Q_k(\phi)=Q_k(\phi)
             +\frac{\|\operatorname{SN}_k\phi\|_2^2}{2k^2}.
\]
The sup-norm and Fubini--Study identities in
\cite[Proposition~35]{Yao25MaximalDestabilizers} give
$\operatorname{FS}_k\operatorname{SN}_k\phi\leq\phi$ and
preserve the section sup-norm after this replacement. Energy
monotonicity therefore reduces minimization to the Fubini--Study
image. On its saturated norm representatives,
\[
 \widetilde Q_k(\operatorname{FS}_k\chi)
       =2M_{X,k}(\chi)+\frac{\|\chi\|_2^2}{2k^2}.
\]
If these embeddings are Chow-unstable, the unique finite-dimensional
minimizer is the maximal Chow-destabilizer with scale
\[
 \|\chi_k^\flat\|_2^2=-2k^2M_{X,k}(\chi_k^\flat).
\]
Saturation ensures that relaxing the image constraint leaves this
minimizer in the required image. Under A1--A4,
\cite[Theorem~4]{Yao25MaximalDestabilizers} concludes convergence
of $\operatorname{FS}_k(\chi_k^\flat)$ to $\phi^\flat$.

A4 is a convergence assumption, not a consequence of the pointwise
expansion in A3. Likewise the convexity proved in
\cite[Proposition~47]{Yao25MaximalDestabilizers} is not by itself
a compactness or convergence theorem for minimizers. The source
does not specify a fixed named topology for A4, so no uniform,
$d_1$ or $d_2$ conclusion is added here. A proof of that passage
must choose the topology and establish the required control of
minimizing sequences. This leaves a precise quantization direction
for further research while retaining the proved finite-dimensional
results.

\subsection{Complex distances and initial velocity}
\label{sec:ytd-complex-distance-regularity}

The complex metric geometry gives a useful way to state a convergence
problem precisely. Let $X$ now be a compact K\"ahler manifold of
dimension $n$ and let $\theta$ be a smooth closed real $(1,1)$-form
whose class is big. Put $V=\operatorname{vol}(\{\theta\})>0$ and
$V_\theta=P_\theta(0)$. Write $\theta_u^n$ for the nonpluripolar
Monge--Amp\`ere measure of $u$. For $1\leq p<\infty$, the complex
finite-energy space is
\[
 \mathcal E^p(X,\theta)=
 \left\{u\in\operatorname{PSH}(X,\theta):
   \int_X\theta_u^n=V,\quad
   \int_X|u-V_\theta|^p\theta_u^n<\infty\right\}.
\]
Here $u$ is an additive complex potential, so its curvature is
$\theta+\ddc u$.

Gupta constructs a complete geodesic metric on this space. We use
the mass-normalized distance $d_p^{\mathrm C}=V^{-1/p}d_p^{\mathrm G}$,
where $d_p^{\mathrm G}$ is the distance in the cited source. Then
\cite[Definition~4.10, Lemma~4.12 and Theorem~5.1]{Gupta24GeodesicV1}
give
\[
 C_{n,p}^{-1}I_p^{\mathrm C}(u,v)
 \leq d_p^{\mathrm C}(u,v)^p
 \leq C_{n,p}I_p^{\mathrm C}(u,v),\qquad
 I_p^{\mathrm C}(u,v)=\frac1V
       \int_X|u-v|^p(\theta_u^n+\theta_v^n).
\]
In particular, uniform convergence of bounded potentials in this space
implies convergence for $d_p^{\mathrm C}$: if
$\sup_X|u_j-u|\leq\delta_j\to0$, the last integral is at most
$2\delta_j^p$.
Theorem~5.4 of the same version says that the weak geodesic $u_t$
joining any $u_0,u_1\in\mathcal E^p(X,\theta)$ satisfies
\[
 d_p^{\mathrm C}(u_s,u_t)=|s-t|d_p^{\mathrm C}(u_0,u_1),
 \qquad s,t\in[0,1].
\]
The distance on general endpoints is defined by decreasing regular
envelopes and a limit of their distances.

The initial-velocity formula has a more specific hypothesis.
For example, if $u_0=P_\theta(f_0)$ with $f_0\in C^2(X)$ and
$u_1\in\mathcal E^p(X,\theta)$, then
\cite[Lemma~5.3]{Gupta24GeodesicV1} gives
\[
 d_p^{\mathrm C}(u_0,u_1)^p
   =\frac1V\int_X|\dot u_{0+}|^p\theta_{u_0}^n.
\]
This is a special case of its regular obstacle-envelope hypothesis.
The source uses contact-set identities and convergence of velocities
with a fixed regular initial endpoint to establish the formula.
Its extension of the distance and metric-geodesic identity to general
finite-energy endpoints does not remove that hypothesis from the
initial-velocity formula. An identification of a limiting spectral
law therefore needs its own control of velocities and their
Monge--Amp\`ere measures.

For $p=2$, \cite[Theorem~8.3]{Gupta24GeodesicV1} gives the stronger
metric inequality
\[
 \begin{split}
 d_2^{\mathrm C}(u,v_t)^2\leq{}&
 (1-t)d_2^{\mathrm C}(u,v_0)^2+t d_2^{\mathrm C}(u,v_1)^2\\
 &-t(1-t)d_2^{\mathrm C}(v_0,v_1)^2,
 \end{split}
\]
where $v_t$ is the weak geodesic and all three endpoints belong to
$\mathcal E^2(X,\theta)$. Thus this complex space is CAT$(0)$.
These results supply complex metric control once a sequence and
its convergence have been established. Applying them to the
quantization problem still requires a comparison with its
non-Archimedean objects and control of the minimizing sequences
in A4. Completeness and the existence of geodesics do not supply
that comparison.

\subsection{Bergman measures and uniform curvature bounds}
\label{sec:ytd-bergman-measure-curvature}

There is a second distinction in the quantization problem: convergence
of a measure on $X$ differs from convergence of a spectral law or a
minimizing direction. The identified preprint
\cite{BlockiDarvas26Radon} gives both a full-mass measure convergence
theorem and a curvature estimate under stronger hypotheses. We state
these known results in the book's curvature and probability conventions.

Let $X$ be a compact complex manifold of dimension $n\geq1$, and let
an ample line bundle $L$ carry a smooth positive metric $h$, with
$\omega=-\ddc\log h$ and $V=\int_X\omega^n$. Thus
$\ddc=\rmi\partial\bar\partial/(2\pi)$ as in the
\hyperref[front:conventions-curvature]{front-matter convention}.
For a full-mass potential
$\varphi\in\mathcal E(X,\omega)$ define
\[
\tag{68.23a}\label{eq:ytd-full-mass-bergman-data}
\begin{split}
 H^k_\varphi(s,s')&=\int_Xh^k(s,s')\rme^{-k\varphi}\omega^n,\\
 K^k_\varphi&=\sum_{j=1}^{N_k}|s_j|_{h^k}^2,
 \qquad P^k_\varphi=k^{-1}\log K^k_\varphi,
 \qquad N_k=h^0(X,L^k),
\end{split}
\]
where $(s_j)$ is $H^k_\varphi$-orthonormal. Full mass gives zero
Lelong numbers, and exponential integrability makes the Hilbert form
finite on the entire section space. We take $k$ sufficiently large
that $L^k$ is very ample. Then $K^k_\varphi$ is a positive smooth
function, and $\omega+\ddc P^k_\varphi$ is a smooth positive
Fubini--Study form. Ampleness alone does not supply this assertion
at every small degree.

The Bergman probability measure is
\[
\tag{68.23b}\label{eq:ytd-bergman-probability-limit}
 \beta^k_\varphi
   =\frac{K^k_\varphi\rme^{-k\varphi}}{N_k}\omega^n,
 \qquad \int_X\beta^k_\varphi=1,
 \qquad
 \beta^k_\varphi\rightharpoonup V^{-1}\operatorname{MA}_\omega(\varphi).
\]
The convergence is the probability form of
\cite[Theorems~1.4 and~5.4]{BlockiDarvas26Radon}; it requires full
mass, but no finite-energy hypothesis. The finite-level normalization
is exact, since each basis section has squared Hilbert norm one.
The density is defined almost everywhere and is integrable. It need
not be smooth when $\varphi$ is singular, even though the kernel and
the Fubini--Study form are smooth. This measure is also different
from the top wedge of that form.

In particular, for a nonpluripolar probability measure $\nu$, the
complex full-mass Calabi--Yau solution
$\operatorname{MA}_\omega(\varphi_\nu)=V\nu$ gives
$\beta^k_{\varphi_\nu}\rightharpoonup\nu$.
The prescribed measure may be singular relative to a smooth volume
form. The conclusion is weak convergence tested against continuous
functions, without a claim of convergence of densities. Arbitrary
K\"ahler classes without a positive line bundle are outside this
construction, and partial kernels for non-full-mass singularities
are not included in this statement.

For the separate regularity result, let $\varphi$ have bounded
Laplacian, denoted $C^{1,\bar1}$ in the source, and assume the current
inequalities $0<a\omega\leq\omega+\ddc\varphi\leq A\omega$.
The known estimate is
\[
\tag{68.23c}\label{eq:ytd-bergman-curvature-bounds}
 \frac{a^{n+1}}{C(1+A^n)}\omega
 \leq\omega+\ddc P^k_\varphi
 \leq\frac{C(1+A^{n+1})}{a^n}\omega,
 \qquad k\geq C/a,
\]
with $C$ depending only on the fixed polarized background
\cite[Theorem~1.1]{BlockiDarvas26Radon}. Its consequence is
\[
\tag{68.23d}\label{eq:ytd-bergman-c1alpha}
 P^k_\varphi\longrightarrow\varphi
 \quad\hbox{in }C^{1,\alpha}(X),\qquad 0<\alpha<1.
\]
Bounded Laplacian is not the assertion that all real second
derivatives are bounded. The conclusion does not include the endpoint
$C^{1,1}$ norm or convergence of curvature tensors.

To compare normalizations, the source's curvature operator is
$\rmi\partial\bar\partial$ and its background is $2\pi\omega$;
the potential $\varphi$ and weight $\rme^{-k\varphi}$ are the same.
Rescaling the Hilbert volume by $(2\pi)^n$ multiplies the kernel by
$(2\pi)^{-n}$, changes $P^k_\varphi$ by a constant of order $k^{-1}$,
and leaves both the curvature and the probability measure unchanged.
Using $N_k$ avoids importing a dimension-dependent raw-volume factor.

The source obtains the curvature estimate from local kernel and
first-jet estimates, and the full-mass convergence from finite-energy
quantization together with a uniform small-mass estimate on deep
sublevel sets. Neither conclusion asserts an initial-velocity
pushforward formula at an arbitrary continuous potential. They give
no compactness or convergence of the minimizing directions in (A4),
and do not prove (A2) or (A3). The positive lower bound $a>0$ is part
of the regularity statement; it cannot simply be dropped for a
degenerate time slice of a weak geodesic. In the identified February
2026 version, decoupling the two curvature bounds and extending
partial-kernel convergence beyond full mass are proposed directions,
rather than conclusions used here.

\section{Adapted measures and big-class stability}
\label{sec:ytd-klt-adapted-stability}

The complex variational bridge extends beyond smooth polarized
manifolds. The identified preprint \cite{DarvasZhang26KltV2},
arXiv:2608.23505v2, revised 13 September 2026, supplies the
continuity input for adapted measures on klt pairs and removes a
semipositivity requirement from a big-class existence criterion.
The results in this section are imports from that version. Their
complex analytic domain is distinct from the general
non-Archimedean envelope input used elsewhere in the book.

\subsection{Local continuity with an adapted measure}

Let $(X,\Delta)$ be a connected normal complex analytic klt pair
of pure dimension $n\geq1$. Here $\Delta\geq0$ is an effective
$\mathbf R$-Weil divisor, $K_X+\Delta$ is $\mathbf R$-Cartier,
and the discrepancies on a log resolution $\pi:Y\to X$ are
$a_i>-1$. A smooth formal Hermitian metric on $K_X+\Delta$
defines an adapted measure $\mu$, with no mass on $X_{\rm sing}$.
Locally on the resolution its pullback has the form
\[
 d\widetilde\mu=\rme^G\prod_i|z_i|^{2a_i}\,dV_Y,
 \qquad \pi_*\widetilde\mu=\mu,
\]
where $G$ is smooth and the divisors $z_i=0$ have simple normal
crossings. The discrepancy inequalities ensure local integrability.

For an open $\Omega\subset X$ and $K\Subset\Omega$, set
\[
 c_{\mu,K}[u]=\sup\{c>0:\rme^{-cu}\in L^1(U,\mu)
       \text{ for some }K\subset U\Subset\Omega\}.
\]
Convergence $u_j\to u$ in $L^1_{\rm loc}(\Omega)$ is measured
using a local K\"ahler volume from an analytic embedding, rather
than being assumed with respect to $\mu$.
By \cite[Theorem~1.1]{DarvasZhang26KltV2}, if
$u_j,u\in\operatorname{PSH}(\Omega)$ converge in this sense and
$\rme^{-\lambda u}$ is integrable on a neighborhood of $K$, then
$\lambda<c_{\mu,K}[u]$. For each
$\lambda<t<c_{\mu,K}[u]$, some $K\subset W\Subset\Omega$
satisfies, for all sufficiently large $j$,
\[
\tag{68.23e}\label{eq:ytd-klt-exponential-continuity}
 \rme^{-tu_j},\rme^{-tu}\in L^1(W,\mu),
 \qquad \rme^{-\lambda u_j}\longrightarrow\rme^{-\lambda u}
           \quad\text{in }L^1(W,\mu).
\]

The proof first obtains a uniform integral bound at the larger
exponent $t$. It slices by regular levels of a holomorphic
function whose zero set contains the singular and boundary-divisor
loci. The classical smooth Demailly--Koll\'ar theorem applies to
these slices. An $L^2$ extension estimate, quantitative division
and normal-family estimates on the normal space give the uniform
bound on a fixed neighborhood. Convergence in measure and this
higher-exponent bound then give the displayed $L^1$ convergence.
A resolution alone does not prove the result: arbitrary weights
$\rme^{\chi-\psi}$ on a smooth resolution do not have the required
continuity. The intrinsic adapted-measure hypotheses are essential
to this proof route.

\subsection{Analytic and divisorial thresholds}

Now assume $X$ is compact normal K\"ahler. Retain the klt pair
and let $\theta$ be a smooth closed real $(1,1)$ form with big
class. Write $V=\operatorname{vol}(\{\theta\})>0$ and
$V_\theta=\sup\{v\in\operatorname{PSH}(X,\theta):v\leq0\}$.
The energy on $\mathcal E^1(X,\theta)$ is normalized by mass:
\[
 I_\theta(u)=\frac1{(n+1)V}\sum_{j=0}^n
 \int_X(u-V_\theta)
 \left\langle\theta_u^j\wedge\theta_{V_\theta}^{n-j}\right\rangle,
 \qquad I_\theta(V_\theta)=0.
\]
Choose a quasi-psh $\psi$ such that
$\mu_\psi=\rme^{-\psi}\mu$ is a probability measure. The analytic
thresholds are
\[
\tag{68.23f}\label{eq:ytd-klt-analytic-thresholds}
\begin{aligned}
 \alpha^A_\psi
 &=\sup\left\{a>0:
   \sup_{\substack{u\in\operatorname{PSH}(X,\theta)\\\sup_Xu=0}}
                \int_X\rme^{-au}\,d\mu_\psi<\infty\right\},\\
 \delta^A_\psi
 &=\sup\left\{d>0:
   \sup_{u\in\mathcal E^1(X,\theta)}
        \int_X\rme^{-d(u-I_\theta(u))}\,d\mu_\psi<\infty\right\}.
\end{aligned}
\]
The second expression is invariant under adding a constant to $u$;
it is an optimal Moser--Trudinger exponent on finite-energy
potentials, rather than on all singularity types.

For a prime divisor $E$ on a proper bimeromorphic model
$\pi:Y\to X$, let
$A_\psi(E)=A_{X,\Delta}(E)-\nu(\psi,E)$, with generic Lelong
numbers normalized so that $\log|z|^2$ has divisorial coefficient
one. Define
\[
\tag{68.23g}\label{eq:ytd-klt-divisorial-thresholds}
\begin{aligned}
 \tau_\theta(E)&=\sup\{s\geq0:\pi^*\{\theta\}-s\{E\}
                                 \text{ is big}\},\\
 S_\theta(E)&=\frac1V\int_0^{\tau_\theta(E)}
           \operatorname{vol}(\pi^*\{\theta\}-s\{E\})\,ds,\\
 \alpha_\psi&=\inf_E\frac{A_\psi(E)}{\tau_\theta(E)},
 \qquad \delta_\psi=\inf_E\frac{A_\psi(E)}{S_\theta(E)}.
\end{aligned}
\]
The infima range over all such prime divisors, without a
dreaminess requirement. An infinite quotient is used when its
denominator vanishes. Theorems~1.2 and~1.3 of
\cite{DarvasZhang26KltV2} identify
\[
             \alpha^A_\psi=\alpha_\psi,
             \qquad \delta^A_\psi=\delta_\psi.
\]
The local continuity theorem is the compactness input for the
uniform analytic estimates. The divisorial formulas concern this
fixed complex klt pair and its big class.

\subsection{The ray argument without convexity of the twist}

For $\lambda>0$ in the integrability range, write
\[
\tag{68.23h}\label{eq:ytd-klt-ding}
 D^\lambda(u)=-\frac1\lambda
        \log\int_X\rme^{-\lambda u}\,d\mu_\psi-I_\theta(u).
\]
It is invariant under additive constants. The radial formula of
\cite[Theorems~1.4 and~4.1]{DarvasZhang24TwistedV4}, transferred
to the adapted measure on a resolution as in
\cite[Section~7.2]{DarvasZhang26KltV2}, reads
\[
\tag{68.23i}\label{eq:ytd-klt-ding-radial}
 \liminf_{t\to\infty}\frac{D^\lambda(u_t)}t
 =-I_\theta\{u_t\}
    +\sup\left\{\tau:\int_X\rme^{-\lambda\widehat u_\tau}
                                         \,d\mu_\psi<\infty\right\},
\]
for a sublinear finite-energy subgeodesic ray from $V_\theta$,
where $\widehat u_\tau=\inf_{t>0}(u_t-t\tau)$ and
$I_\theta\{u_t\}=\lim_{t\to\infty}I_\theta(u_t)/t$.
Here $\lambda<c_{\mu_\psi}[V_\theta]$ is retained. This
radial formula and the divisorial semistability threshold do not
require convexity of $D^\lambda$ along rays.

The new step is \cite[Theorem~7.3]{DarvasZhang26KltV2}.
If $D^\lambda$ were unbounded below at
$0<\lambda<\delta_\psi$, normalized bad endpoints would give
weak geodesic segments of lengths $-I_\theta(u_j)\to\infty$.
Their local $L^1$ limit is initially a subgeodesic ray. Proposition~7.2
controls its Legendre transforms at every fixed $\tau<0$.
The adapted-measure continuity then forces
$\int_X\rme^{-\lambda\widehat u_\tau}\,d\mu_\psi=\infty$
for $-1<\tau<0$. Thus the limit cannot be the constant ray.
Maximalizing its test curve preserves the integrability thresholds
and radial energy, and gives a nontrivial finite-energy geodesic ray.
The source obtains negative radial Ding slope at every
$\eta\in(\lambda,\delta_\psi)$, contradicting the known
nonnegative-slope characterization of $\delta_\psi$.
This establishes boundedness below throughout the strict range.

There is also a quantitative H\"older step. For
$\eta>\lambda>0$ and $0<a<\min(\lambda,\alpha_\psi)$,
\cite[Lemma~7.4]{DarvasZhang26KltV2} gives, when $\sup_Xu=0$,
\[
\tag{68.23j}\label{eq:ytd-klt-holder-coercivity}
 D^\lambda(u)\geq
 \frac{\eta(\lambda-a)}{\lambda(\eta-a)}D^\eta(u)
 -\frac{a(\eta-\lambda)}{\lambda(\eta-a)}I_\theta(u)-C.
\]
Choosing $\lambda<\eta<\delta_\psi$ makes $D^\eta$ bounded
below, and $-I_\theta(u)=d_1(V_\theta,u)$ on this normalized
slice. Hence $D^\lambda$ is coercive. The exponential part is
$L^1$ continuous on finite-energy potentials by the adapted
continuity theorem; the variational argument yields a minimizer.
Its Euler--Lagrange equation is
\[
\tag{68.23k}\label{eq:ytd-klt-ke-equation}
 \frac{\left\langle\theta_u^n\right\rangle}{V}
 =\frac{\rme^{-\lambda u}\,d\mu_\psi}
        {\int_X\rme^{-\lambda u}\,d\mu_\psi}.
\]
An additive constant makes this
$\left\langle\theta_u^n\right\rangle
 =\rme^{-\lambda u-\psi}\,d\mu$.
Corollary~7.5 of \cite{DarvasZhang26KltV2} gives minimal
singularities, using the big-class $L^{1+\varepsilon}$ estimate
of \cite[Theorem~B]{BEGZ10BigV2} on a resolution.
Minimal singularities do not mean boundedness across the polar
locus of a merely big class.

In particular, $\delta_\psi>1$ is a sufficient divisorial
criterion at $\lambda=1$, without semipositivity of the twisting
current. Neither existence at the endpoint $\lambda=\delta_\psi$
nor a general necessity or uniqueness statement follows from this
strict-range argument. It supplies a complex klt existence result;
it does not repair the general Siu or envelope assertions, an
initial spectral pushforward, or the minimizing-direction limit
in (A4).

\section{Stability and small-parameter uniqueness in big classes}
\label{sec:ytd-big-small-uniqueness}

The existence criterion just discussed leaves uniqueness as a separate
question. On a smooth compact K\"ahler manifold, the September preprint
\cite{DangZhangZhou26UniquenessV1}, arXiv:2609.00810v1,
gives a sufficiently small parameter range. Its useful analytic input
is an estimate for the difference of two potentials, even when neither
potential is globally bounded. This section imports that estimate and
explains its uniqueness consequence with a weighted integrability
argument. It also records the distinction between power and logarithmic
moduli of continuity.

\subsection{Stability relative to the minimal singularity type}
\label{sec:ytd-big-density-stability}

Fix a connected smooth compact K\"ahler manifold $(X,\omega)$ of
dimension $n\geq1$, a smooth closed real $(1,1)$ form $\theta$ with
big class, and a smooth probability volume $d\nu$. Put
$V=\operatorname{vol}(\{\theta\})>0$ and use the probability
Monge--Amp\`ere operator
\[
 \operatorname{MA}_\theta(u)
       =V^{-1}\left\langle(\theta+\ddc u)^n\right\rangle.
\]
All geometry and the volume form remain fixed in the following
estimates. If $u,v$ have minimal singularities, their difference norm
means the least $S\geq0$ such that $u\leq v+S$ and $v\leq u+S$.
It is finite; this definition avoids subtracting two infinite values
on a common polar set.

For $q>1$, let $F,G\geq0$ be probability densities relative to
$\nu$, with both $L^q(\nu)$ norms at most $B$.
Let $u,v$ have minimal singularities, satisfy $\sup_Xu=\sup_Xv=0$,
and solve $\operatorname{MA}_\theta(u)=F\,d\nu$ and
$\operatorname{MA}_\theta(v)=G\,d\nu$.
The $L^q$ case of \cite[Theorem~3.2]{DangZhangZhou26UniquenessV1}
gives a constant $C_q=C(X,\omega,\theta,\nu,n,q,B)$ such that
\[
\tag{68.23l}\label{eq:ytd-big-density-stability}
 \|u-v\|_\infty
 \leq C_q\left\|F^{1/n}-G^{1/n}\right\|_{L^{nq}(\nu)}.
\]
The source sets its class volume equal to one. The rescaling to the
book's operator and $\ddc$ convention is given in
\Cref{sec:provenance-big-small-uniqueness}; it changes fixed constants.
No uniformity as the big class approaches zero volume is asserted.

The proof uses the extrapolated envelope
$P_\theta(bu-(b-1)v)$ for $b>1$. Since both potentials are
$V_\theta$ plus bounded functions, its obstacle is also
$V_\theta$ plus a bounded function. The envelope therefore has a
psh minorant and minimal singularities. Contact concentration and
the mixed Monge--Amp\`ere inequality bound its density in terms of
$b|F^{1/n}-G^{1/n}|$ and $G^{1/n}$. The big-class
$L^q$ estimate then controls its oscillation relative to
$V_\theta$, and full mass forces a useful contact point.
This yields the displayed difference estimate. The contact and
normalization inputs can be checked in
\cite[Theorem~2.7, Lemma~2.9 and Lemma~3.5]{DDL23RelativeV2};
these are locators in the acquired arXiv version, not substituted
numbers from the published survey.

\subsection{A weighted contraction for the mean-field equation}
\label{sec:ytd-big-weighted-contraction}

Let $F\geq0$ be a probability density in $L^p(\nu)$ for some
$p>1$. Theorem~3.3 of \cite{DangZhangZhou26UniquenessV1}
asserts existence and uniqueness with minimal singularities for
\[
\tag{68.23m}\label{eq:ytd-big-small-mean-field}
 \operatorname{MA}_\theta(u)=\rme^{-\lambda u}F\,d\nu,
 \qquad 0<\lambda<\lambda_0,
\]
where $\lambda_0>0$ depends on the fixed geometry, $p$ and a
bound for $\|F\|_{L^p(\nu)}$. The parameter is sufficiently
small, rather than the whole strict existence range from a
divisorial criterion. Here the equation itself fixes the additive
constant.

The following calculation is an explicitly identified direct
consequence of \eqref{eq:ytd-big-density-stability} and uniform
Skoda integrability. It provides the uniqueness step without a
global supremum of exponentials of unbounded potentials.
Choose $1<q<p$ and put $r'=p/(p-q)$. Uniform integrability on
the normalized psh slice supplies $\alpha>0$ and $C_\alpha$ with
\[
\tag{68.23n}\label{eq:ytd-big-uniform-exponential}
 \int_X\rme^{-\alpha w}\,d\nu\leq C_\alpha,
 \qquad w\in\operatorname{PSH}(X,\theta),\quad\sup_Xw=0.
\]
This is the fixed-geometry form of
\cite[Theorem~2.2]{DangZhangZhou26UniquenessV1}.
For $q\lambda r'\leq\alpha$, H\"older's inequality gives
$\|\rme^{-\lambda w}F\|_{L^q(\nu)}
 \leq\|F\|_{L^p(\nu)}C_\alpha^{1/(qr')}$,
uniformly in $w$ and in this parameter interval.

Suppose $u,v$ solve \eqref{eq:ytd-big-small-mean-field} with
minimal singularities. Set $u_0=u-\sup_Xu$, $v_0=v-\sup_Xv$,
and $S=\|u_0-v_0\|_\infty$. Define
\[
\tag{68.23o}\label{eq:ytd-big-normalized-densities}
 Z_w=\int_X\rme^{-\lambda w}F\,d\nu,
 \qquad H_w=Z_w^{-1}\rme^{-\lambda w}F,
 \qquad w\in\{u_0,v_0\}.
\]
Since $w\leq0$ and $\int_XF\,d\nu=1$, one has $Z_w\geq1$.
The preceding estimate bounds both probability densities $H_w$
in $L^q(\nu)$ by a fixed $B\geq1$. Moreover the comparison
$|u_0-v_0|\leq S$ away from their polar set implies
$|\log Z_{u_0}-\log Z_{v_0}|\leq\lambda S$.
The elementary exponential difference inequality therefore gives,
almost everywhere,
\[
 |H_{u_0}^{1/n}-H_{v_0}^{1/n}|
 \leq\frac{2\lambda S}{n}
                 (H_{u_0}^{1/n}+H_{v_0}^{1/n}).
\]
The polar set has zero smooth volume; the comparison extends to
the potentials by their psh representatives. Taking the $L^{nq}$
norm and applying \eqref{eq:ytd-big-density-stability} yields
\[
\tag{68.23p}\label{eq:ytd-big-weighted-contraction}
 S\leq\frac{4C_qB^{1/n}}n\lambda S.
\]
Choose $\lambda_0$ smaller than both $\alpha/(qr')$ and
$n/(4C_qB^{1/n})$. Then $S=0$. Equality of the normalizing
integrals and the unit mass condition give
$\lambda\sup_Xu=\log Z_{u_0}=\log Z_{v_0}=\lambda\sup_Xv$,
so $u=v$. Each constant used here is uniform over the two
solutions. Minimal singularities ensure a finite difference norm;
they do not provide a finite supremum of $\rme^{-\lambda u_0}$.

Existence can likewise be recovered as a direct formal consequence
of these imported inputs and the full-mass solvability theorem
\cite[Theorem~A]{BEGZ10BigV2}. Use the slice
\[
 \mathcal K=\left\{w\in\operatorname{PSH}(X,\theta):
                            \int_Xw\,d\nu=0\right\}.
\]
It is convex and compact in $L^1(\nu)$, and its elements have a
uniformly bounded supremum by normalized psh compactness.
The formula for $Z_w$ and $H_w$ is used also for $w\in\mathcal K$.
The slice defined by $\sup_Xw=0$ generally is not convex;
integral normalization is essential at this step.
Choose $1<q<q_1<p$ and decrease the parameter bound so that
H\"older and \eqref{eq:ytd-big-uniform-exponential} give uniform
$L^{q_1}$ bounds for $\rme^{-\lambda w}F$ on $\mathcal K$.
The uniform upper bound for $w$ also gives a positive lower bound
for $Z_w$. Thus the probability densities $H_w$ have uniform
$L^{q_1}$ bounds. If $w_j\to w$ in $L^1$, almost-everywhere
subsequences and this higher integrability give
$H_{w_j}\to H_w$ in $L^q$.

Let $T(w)$ be the unique integral-normalized full-mass solution
of $\operatorname{MA}_\theta(T(w))=H_w\,d\nu$.
The $L^{q_1}$ density makes it minimally singular by
\cite[Theorem~B]{BEGZ10BigV2}. The elementary bound
$|a^{1/n}-b^{1/n}|^n\leq|a-b|$ for $a,b\geq0$ and
\eqref{eq:ytd-big-density-stability} give convergence of the
supremum-normalized output potentials in their difference norm.
Integral normalization preserves convergence, so
$T:\mathcal K\to\mathcal K$ is continuous. Schauder's theorem
gives $w=T(w)$, and $u=w+\lambda^{-1}\log Z_w$ solves
\eqref{eq:ytd-big-small-mean-field}. The fixed point has minimal
singularities because it belongs to the image of $T$.
This calculation changes the normalization in the source's
fixed-point paragraph and avoids its auxiliary positive-exponent
stability proof. The resulting density $H_{u_0}$
is in $L^q$. Consequently
\cite[Theorem~D]{DDGHKZ14HolderPublished} gives local H\"older
continuity on $\operatorname{Amp}(\{\theta\})$, with every
exponent $0<\beta<2/(1+nq')$, where $q'=q/(q-1)$.
This conclusion concerns the ample locus of a smooth manifold;
it asserts neither global boundedness in a merely big class nor
uniqueness for arbitrary normal klt spaces or all positive parameters.

\subsection{Logarithmic continuity under weaker density assumptions}
\label{sec:ytd-big-log-modulus}

There is also a useful weaker regularity statement when a power
$L^p$ bound is unavailable. Let $F\geq0$ have unit integral and
\[
\tag{68.23q}\label{eq:ytd-big-log-density}
 \int_X F[\log(1+F)]^a\,d\nu<\infty,
 \qquad a>n.
\]
Suppose $\operatorname{MA}_\theta(\varphi)=F\,d\nu$,
$\sup_X\varphi=0$, and retain the explicit minimal-singularity
bound $V_\theta-M_0\leq\varphi\leq V_\theta$.
Theorem~4.1 and the proof of Theorem~1.2 in
\cite{DangZhangZhou26UniquenessV1} give, for every compact
$K\Subset\operatorname{Amp}(\{\theta\})$ and
$0<\gamma<a/n-1$, constants $C_K$ and $0<d_K<1$ such that
\[
\tag{68.23r}\label{eq:ytd-big-log-modulus}
 |\varphi(x)-\varphi(y)|
 \leq\frac{C_K}{|\log d_\omega(x,y)|^\gamma},
 \qquad x,y\in K,\quad0<d_\omega(x,y)<d_K.
\]
The constants may depend on $M_0$, the entropy integral and the
fixed geometry. This formulation retains minimal singularities
as a hypothesis; it does not use an unacquired auxiliary preprint
to deduce that bound from the entropy integral.

The source first bounds $\sup_X(\psi-\varphi)$ in terms of the
$L^1$ size of $(\psi-\varphi)_+$ for normalized competitors.
An auxiliary equation, uniform exponential integrability and
an entropy inequality make the Monge--Amp\`ere mass of sublevel
sets vanish at a controlled level. Domination converts that mass
estimate to a potential estimate. Demailly--Kiselman regularization
then supplies a logarithmic bound for $\rho_\delta\varphi-\varphi$
on a slightly larger compact subset of the ample locus.
In this construction the strictly positive analytic barrier may
first be shifted down so that it lies below every relevant
regularized obstacle; this preserves its curvature and justifies
the comparison of positive-part $L^1$ errors.

The passage from averages to oscillations uses nested interior
sets and the submean argument in
\cite[Lemma~4.4]{LPT21StabilityPublished} and
\cite[Lemma~3.3]{Zeriahi20ModulusV1}. For
$\kappa(t)=|\log t|^{-\gamma}$ the growth condition (3.1) in
the latter source holds after restricting to small $t$.
Its local lemma avoids a boundary regularity assertion.
The logarithmic estimate alone does not yield H\"older continuity;
the earlier power modulus uses the stronger $L^q$ density.
None of these estimates identifies an initial geodesic velocity
with a spectral measure, controls a limit of destabilizing
minimizers, or supplies the general non-Archimedean Siu input.

\section{Dated 2026 claims}
\label{sec:ytd-recent-claims}

The following entries report source versions available at the fixed
22 August 2026 cutoff. They are not updated silently when a preprint changes
version, and they do not enlarge the theorem domains in the preceding
sections.

\begin{remark}[Boucksom--Jonsson, version 2]
\label{rem:ytd-bj26}
The preprint \cite{BJ26ComplexNA}, version~2 dated 19 May 2026,
develops a broad complex-to-non-Archimedean approach to the YTD
conjecture. In this manuscript it is used for orientation and
terminology only. Its later bridge claims have not been substituted
for the dependency-audited theorems above.
\end{remark}

\begin{remark}[Trusiani, version 2]
\label{rem:ytd-trusiani26}
The preprint \cite{Trusiani26}, version~2 dated 17 August 2026, retains its
automorphism-relative uniform YTD theorem for smooth polarized
projective varieties and an entropy-regularization theorem; its
Corollary~B records the equivalence of the smooth projective stability
signatures stated there. Its Corollary~A gives entropy regularization;
the following discussion derives the projective entropy-slope equality
for maximal finite-energy rays. The new
appendix by S\'ebastien Boucksom gives a non-Archimedean proof of entropy
regularization over $\Bbbk((t))$, while Theorem~C adds the
weighted-relative correspondence and the log-concave weighted-extremal
consequence.  This book records those statements as dated recent claims. We
have not audited every dependency or promoted them to timeless theorems, and
they do not by themselves treat an arbitrary nonprojective compact K\"ahler
class or identify separately prescribed canonical approximation sequences
term by term.
\end{remark}

\begin{proposition}\label{prop:ytd-status}
The status used in this book is:
\begin{enumerate}
\item $E,J$-slopes are theorems for normal ample test
      configurations.
\item All smooth-ray slopes are theorems under positivity alone, by
      \cite[Theorem~3.6]{BHJ19}; the subklt hypothesis of
      \cite[Theorem~4.2]{BHJ19} is needed only for the logarithmic
      extension.
\item The unrelative Fano Ding equivalence is a theorem for smooth
      projective Fano varieties with finite automorphism group; the
      semipositive twisted/logarithmic source has its own stated scope.
\item Model stability implying cscK is a 2025 theorem for all big
      models in its unrelative coercive scope.
\item Automorphism-relative projective uniform YTD, projective entropy regularization and
      the resulting entropy-slope equality, and the weighted-relative
      projective correspondence are recorded as Trusiani's dated 2026
      version-two claims, not as audited theorems here.
\item Arbitrary compact-K\"ahler ordinary YTD is not asserted.
\end{enumerate}
\par
\end{proposition}

\begin{proof}
The first row is \Cref{thm:ytd-energy-slopes}; the second is
\Cref{thm:ytd-smooth-functional-slopes}; the third is
\Cref{thm:ytd-fano-bridge}; and the fourth is
\Cref{thm:ytd-mpwn-implication}. The final two rows are the
scope statements of \Cref{rem:ytd-bj26,rem:ytd-trusiani26}.
\end{proof}

\section{A bridge diagram}
\label{sec:ytd-diagram}

The audited implications can be displayed without hiding their
hypotheses:
\[
\begin{array}{ccccc}
\text{big model data}
 &\longrightarrow&
\mathcal E^{1,\NA}(\alpha)
 &\xrightarrow{\operatorname{MA}^{\NA}}&
\mathcal M^{1,\NA}(\alpha)\\
\big\downarrow
 &&
\big\downarrow{\scriptstyle M^{\NA},J^{\NA}}
 &&
\big\downarrow{\scriptstyle\beta_\alpha,E_\alpha^\vee}\\
\text{DF and model norm}
 &\longrightarrow&
\text{finite-energy coercivity}
 &\longrightarrow&
\text{cscK metric}.
\end{array}
\tag{68.24}\label{eq:ytd-audited-diagram}
\]
The left-to-middle completion requires \emph{big} model testing.
Replacing the upper-left corner by ordinary ample test configurations
breaks the proved implication.

\begin{example}[Product configuration]
\label{ex:ytd-product-configuration}
A one-parameter subgroup of $\Aut(X,L)$ gives a product test
configuration. Its non-Archimedean $J$-norm can be positive before
reduction by the automorphism group, while the Mabuchi functional is
affine along the induced orbit. This is why a stability statement with
continuous automorphisms must use a reduced norm and a relative
functional.

On $(\mathbf P^1,\mathcal O(1))$, the action
$[x:y]\mapsto[\tau x:\tau^{-1}y]$ gives weights $2k-m$ on the
$m+1$ degree-$m$ monomials. Their average is zero and their normalized
maximum is one. Thus $E^{\NA}=0$ and $J^{\NA}=1$; the limiting weight
measure is uniform on $[-1,1]$.

The same orbit gives a useful analytic normalization check. Write
$u=|x|^2/(|x|^2+|y|^2)$; normalized Fubini--Study measure sends $u$ to
Lebesgue measure on $[0,1]$. A squared potential is
\[
 \varphi_s=\log\bigl(\rme^{2s}u+\rme^{-2s}(1-u)\bigr).
\]
After change of variables by the automorphism, its energy derivative is
$2\int_0^1(2u-1)\,\rd u=0$, so $E_\omega(\varphi_s)=0$.
For $s>0$, direct integration gives
\[
 J_\omega(\varphi_s)=2s\coth(2s)-1,
 \qquad \frac{J_\omega(\varphi_s)}{2s}\longrightarrow1.
\tag{68.24a}\label{eq:ytd-product-orbit-j}
\]
Its density is $q_s(u)=(\rme^{2s}u+\rme^{-2s}(1-u))^{-2}$.
Substitution in the entropy integral gives
$H_\omega(\varphi_s)=2s\coth(2s)-1$. Since
$\operatorname{Ric}\omega=2\omega$, the Ricci formula gives
$R_\omega(\varphi_s)=-J_\omega(\varphi_s)$. Hence $M_\omega=0$ along
this nonconstant orbit, as expected for cscK metrics. Full relative entropy
in place of $H_\omega$ would fail this cancellation. Reducing by the
automorphism orbit makes the reduced $J$ vanish.
\end{example}

\begin{example}[Nonreduced central fibre]
\label{ex:ytd-nonreduced-central}
Normalized ramified base change can remove existing multiplicity. At
the generic point of a component with $t$ of order $b$, degree $d$ gives
order $b/\gcd(b,d)$ on each normalized branch; in particular a reduced
component remains reduced. This explains the sufficiently divisible
degree in \Cref{prop:ytd-base-change-correction}.

For an explicit nonreduced example, take the quotient
\[
 \mathcal X=(\mathbf P^1_z\times\mathbf A^1_u)/\boldsymbol\mu_m,
 \qquad \zeta\cdot(z,u)=(\zeta z,\zeta u),\qquad t=u^m,
 \quad m\geq2.
\]
The pullback of $\mathcal O_{\mathbf P^1}(m)$ descends to a relatively
ample line bundle $\mathcal L$: the stabilizers at $z=0,\infty$ act
trivially on its fibres. The quotient is normal and flat over the curve
base; its general polarized fibre is $(\mathbf P^1,\mathcal O(m))$.
Diagonal scaling by an $m$th root of $\lambda$ descends to a well-defined
test-configuration action covering $t\mapsto\lambda t$.
At a general point of $u=0$ the quotient has no stabilizer, so
$\mathcal X_0=mG$. The degree-$m$ map $\mathbf P^1\to G$ pulls
$\mathcal L|_G$ back to $\mathcal O(m)$; hence
$(G\cdot\mathcal L)=1$, $V=m$ and
\[
 C(\phi)=\frac{1-m}{m}.
\tag{68.24b}\label{eq:ytd-quotient-correction}
\]
The normalized degree-$m$ base change is the product configuration on
$(\mathbf P^1,\mathcal O(m))$ with action $z\mapsto\lambda z$.
Its weights $0,1,\ldots,mk$ in degree $k$ give $J^{\NA}=m/2$;
its Mabuchi invariant is zero by \Cref{thm:ytd-smooth-functional-slopes},
because the Fubini--Study orbit is cscK and has constant Mabuchi value.
Homogeneity therefore gives, downstairs,
\[
 J^{\NA}(\phi)=\frac12,\qquad M^{\NA}(\phi)=0,
 \qquad \DF(\phi)=\frac{m-1}{m}.
\tag{68.24c}\label{eq:ytd-quotient-df}
\]
Thus $\DF(\phi_m)=0$ while $m\DF(\phi)=m-1$. The correction restores
the homogeneous Mabuchi value. This example also keeps its ordinary
automorphism norm separate from the reduced norm.
\end{example}

\section*{Source notes}
\label{sec:ytd-source-notes}

The finite-energy geodesic convexity used above is
\cite[Theorem~1.2]{BDL17}.  The algebraic and model functionals are
\cite[Definitions~7.13/7.20 and Propositions~7.14--7.22]{BHJ17}.
The analytic entropy is half the ordinary entropy by
\cite[equation~(1.10)]{BHJ19}; $E,J$ double when norm weights become
squared potentials, while the half-entropy and Mabuchi value do not.
Only
the unaffected slope range certified by \cite{BHJ19err} is used.
The Fano variational theorem and its automorphism refinement are
\cite[Theorem~A and Corollary~5.2]{BBJ21}. The transcendental
model, finite-energy, and divisorial-measure bridges are the exact
results of \cite{MP24,MPWN25} already audited in
\Cref{chap:mpwn-k-stability}. The energy formula used in the
valuative construction also uses \cite{DXZ23} on its proved
$\mathcal I$-maximal finite-energy locus. The two 2026 preprints
\cite{BJ26ComplexNA,Trusiani26} remain explicitly dated context at the
22 August 2026 cutoff. In particular, no exact identification of a
canonical entropy approximation or every divisorial-threshold normalization
is inferred from the broader existence claim.
